\subsection*{A uniform pulse-entrance envelope from the prescribed source stages}
Write $P=P_*$, $T=T_d=e^{M_d}+10>10$ and $f=(1+\eta^2)^{-1}$. Retain the normalized moment notation
$a=U/E$, $m=M/(XE)$, $j=J/(XHE)$,
$s=S/(XE^2)$ and $d=(2XE^2)^{-1}\int_X^\infty E^2\,dx$.
These satisfy the exact four transport equations
\[
m'=a-(\ell+\tfrac12)m,\quad j'=a-(2\ell+\tfrac12)j,\quad
s'=a^2-\tfrac12-2\ell s,\quad d'=-\tfrac12-2\ell d.
\]
The prime is differentiation in $\log(X/X_R)$.

\begin{proposition}[Source-native pre-pulse moment bound]\label{prop:uniform-source-prepulse}
For the Appendix-A reference stages (A.7)--(A.9), $P>e^T$ and
$0<\lambda\le1/20$ imply
\[
\sup_{|\eta|\le1}\max\{|m(X_p,\eta)|,|j(X_p,\eta)|\}<\frac9{100}.
\]
\end{proposition}
\begin{proof}
The exact reference integrals at $X_R$ yield
$|m|\le8/P$, $|j|\le5/P$. On the first unit interval,
$U=4\eta$, $0\le\ell\le3/5$, and
$E\ge Pf e^{-y/2}$; consequently $|a|<16/P$.
Both moment equations have damping at least $1/2$, giving
$|m|,|j|\le40/P$ at the end.
During the following axial interval of length $T$,
$0\le U/\eta\le4$, $\ell=0$, and
$|a(y)|\le8e^{(1+y)/2}/P$. Integrating the two equations gives
\[
\max(|m|,|j|)
\le\frac{40e^{-T/2}}P+
\frac{8e^{1/2-T/2}}P(e^T-1)
<40e^{-3T/2}+8e^{(1-T)/2}.
\]
Since $T>10$, the last expression is bounded by
$40e^{-15}+8e^{-9/2}<9/100$.
The final inequality uses exact rational exponential-series lower
bounds $e^{15}>10^6$, $e^{9/2}>90$.
Thereafter $U=0$, while $-\lambda\le\ell\le0$ through the
one-unit transition and intermediate power-law interval.
Both damping coefficients remain positive, so neither absolute
moment increases.
\end{proof}

\begin{proposition}[Explicit smallness of the forward $S$-moment]
\label{prop:source-forward-s-explicit}
Let $T_f>0$ be the source's angular interpolation length and
$k=T+T_f+1$. If
\[
0<\lambda\le\frac{10^{-6}}k,
\]
then $\lambda\sup_{|\eta|\le1}|s(X_p,\eta)|<10^{-4}$.
\end{proposition}
\begin{proof}
At $X_R$, $|s|\le5/12+64/P^2$. The first unit interval has
nonnegative damping $2\ell$ and $|a|<16/P$, giving
$|s|<11/12+320/P^2<1$. On the axial interval,
$\int_0^T a^2\,dy\le64e(e^T-1)/P^2<64e^{1-T}<1$,
hence $|s|<T/2+2$. On the one-unit slope transition,
$|s|<e^{2\lambda}(T/2+5/2)<(9/8)(T/2+5/2)$.
For the subsequent constant slope $-\lambda$ lasting
$W=60\log(1/\lambda)$, one has exactly
\[
s(X_p)=q s_{\rm in}-\frac{q-1}{4\lambda},\qquad
q=e^{120\lambda\log(1/\lambda)}.
\]
For $k\ge11$, monotonicity and $e^{17}>11\cdot10^6$ imply
$120\lambda\log(1/\lambda)<1/5000$, so $q<5000/4999$.
Also $(T+5)/k\le15/11$. Therefore
\[
\lambda|s(X_p)|
<10^{-6}\frac{5000}{4999}\frac9{16}\frac{15}{11}
+\frac1{4\cdot4999}<10^{-4}.
\]
\end{proof}

\begin{proposition}[Conditional backward exterior envelope and A.19 root]
\label{prop:source-full-bias-explicit}
Suppose the actual azimuthal moment corrections meet the source's
stated post-pulse slope bounds: $\ell\le0$ on the angular
interpolation and its following $30\log(1/\lambda)$ plateau,
$\ell$ decreases to $-1$ over one further unit interval,
then equals $-1$ for $4\log(1/h)$, and satisfies
$\ell\le-3h/4$ afterward. For $0<h\le1/100$,
\[
0\le d(X_{\rm end})
\le\frac{T_f}{2}+15\log(1/\lambda)+\frac34+\frac{h^7}{3}.
\]
If also $\lambda\le10^{-6}/(T+T_f+1)$, then
\[
\lambda\sup_{|\eta|\le1}
|s(X_p,\eta)-e^{-26}d(X_{\rm end},\eta)|<\frac1{8000}.
\]
Together with the corrected fixed-$\sigma'$ bump estimate and
Proposition~\ref{prop:uniform-source-prepulse}, the true A.19
remainder satisfies $|\mathcal E|<1/2500$ and
$|\partial_A\mathcal E|<1/100$ on the proposed amplitude interval.
Thus the genuine moment equation has a unique root
$1.009<A(\eta)<1.011$ under these simultaneous stage hypotheses.
\end{proposition}
\begin{proof}
Write $R(v)=XE^2/(X_{\rm end}E_{\rm end}^2)$ for
$v=\log(X/X_{\rm end})$. Then $R(0)=1$, $R'=2\ell R$ and
$d(X_{\rm end})=\tfrac12\int_0^\infty R(v)\,dv$.
The first stages contribute at most
$T_f/2+15\log(1/\lambda)+1/2$; the $-1$ hold contributes at most
$1/4$. At its end $R\le h^8$, and the remaining
$\ell\le-3h/4$ tail contributes at most $h^7/3$.
Set $c=10^{-6}$ and $k=T+T_f+1\ge11$.
The elementary inequalities $\lambda T_f\le c$ and
$\lambda\log(1/\lambda)\le17c/11$ yield
\[
\lambda e^{-26}d(X_{\rm end})\le\lambda d(X_{\rm end})
<\frac c2+\frac{255c}{11}
+\frac c{11}\left(\frac34+\frac1{3\cdot100^7}\right)
<\frac1{40000}.
\]
Together with the preceding forward bound this gives
$1/10000+1/40000=1/8000$. The source-specific corrected
bump contribution is $<5\cdot10^{-17}$ and its amplitude
derivative $<10^{-16}$, because $|m|,|j|<9/100<1$.
Apply the previously established A.19 endpoint signs and
strict derivative inequality.
\end{proof}

The estimate is a theorem about stage bounds, not a fully
frozen numerical source realization. In particular, compatibility
of the choice of $\lambda$ with the source's later selection of $T_f$
and the corrected azimuthal moment bumps must be verified. The
exterior pressure mass, Appendix-B nonlinear contraction,
unforced $n=1,2$ corrections and physical Navier--Stokes residual
remain unverified.
